\newpage
\section{结构}

本宏包定义了两个宏：\tkzMacro{tkz-elements}{tkzGetNodes} 和 \tkzMacro{tkz-elements}{tkzUseLua}。

此外，本宏包还会加载文件 |tkz_elements_main.lua|。该文件初始化了各个模块所使用的全部表，而类正是在这些模块中定义的。它还定义了函数 \tkzFct{tkz-elements}{init\_elements()}，用于重置所有表。


\begin{tikzpicture}[scale=.75]
\begin{scope}
\path[mindmap, concept color=MidnightBlue!60, text=white,text width=38mm,
 level 1 concept/.append style={level distance=120mm,
  sibling angle=72},
 set angles for level/.style={level 2/.append style={
 sibling angle=360/\the\tikznumberofchildren}},
 level/.append style={set angles for level=2},
 level 3 concept/.append style={level distance=20mm,
  sibling angle=20},
 L1/.style={level distance=45mm},
 L2/.style={level distance=65mm,minimum size=2cm}]

node[concept,circular drop shadow] {|tkz-elements.sty|} [clockwise from=10]

child[concept color= MidnightBlue!80,minimum size=4cm,text width=38mm,
clockwise from=27] {
  node[concept,circular drop shadow] {|tkz\_elements\_main|}
  [clockwise from=0]
  child[L2] { node[concept,circular drop shadow] {|Regular \_Polygon|} }
  child[L2] { node[concept,circular drop shadow] {|Quadri\-lateral|} }
  child[L2] { node[concept,circular drop shadow] {|Para\-llelogram|} }
  child[L2] { node[concept,circular drop shadow] {|Rectangle|} }
  child[L2] { node[concept,circular drop shadow] {|Matrix|} }
  child[L2] { node[concept,circular drop shadow] {|Vector|} }
  child[L2] { node[concept,circular drop shadow] {|Point|} }
  child[L2] { node[concept,circular drop shadow] {|Line|} }
  child[L2] { node[concept,circular drop shadow] {|Circle|} }
  child[L2] { node[concept,circular drop shadow] {|Triangle|} }
  child[L2] { node[concept,circular drop shadow] {|Conic|} }
  child[L2] { node[concept,circular drop shadow] {|Square|} }
  child[L2] { node[concept,circular drop shadow] {|OCCS|} }
  child[L2] { node[concept,circular drop shadow] {|Path|} }
};
\end{scope}
\end{tikzpicture}

目前的类有：

 \tkzClass{point} (z)、\tkzClass{circle} (C)、\tkzClass{line} (L)、\tkzClass{matrix} (M)、\tkzClass{occs} (O)、\tkzClass{parallelogram} (P)、\tkzClass{quadrilateral} (Q)、\tkzClass{rectangle} (R)、\tkzClass{square} (S)、\tkzClass{triangle} (T)、\tkzClass{vector} (V)、\tkzClass{conic} (CO)、\tkzClass{regular\_polygon} (RP) 和 \tkzClass{path} (PA)。\\

 括号中的变量表示通常赋予该类对象所在表的名称。

若 \tkzname{name} 指代某个类，则它的定义可在文件 \code{tkz\_elements\_name.lua} 中找到。
\endinput
